% theme: excretion
\MajorQuestion{排出}
\SetRowSpaceQanda

\Instruction{排出のしくみに関する次の問いに答えなさい。}
\QQ{細胞の活動によって体内にできた不要な物質を、体外に出すはたらきを\NamedBlank{excretion}という。}{排出}
\QQ{細胞でできる有害な不要物質のうち、二酸化炭素以外のものを何というか。}{アンモニア}
\QQ{アンモニアは、細胞から何の中に出されるか。}{組織液}
\QQ{アンモニアは、何によって肝臓に運ばれるか。}{血液}
\QQ{アンモニアは、肝臓で毒性の少ない\NamedBlank{urea}に変えられる。}{尿素}

\Instruction{腎臓と尿の通り道に関する次の問いに答えなさい。}
\QQ{腹部の背中側に左右1対ある器官を\NamedBlank{kidney}という。}{腎臓}
\QQ{\RefBlank{kidney}は、血液中から尿素などの不要な物質と余分な水分や塩分をこし出し、何をつくるか。}{尿}
\QQ{\RefBlank{kidney}とぼうこうをつなぐ管を\NamedBlank{ureter}という。}{輸尿管}
\QQ{\RefBlank{ureter}を通って送られた尿を、一時的にためておく器官を\NamedBlank{bladder}という。}{ぼうこう}
\QQ{\RefBlank{bladder}にたまった尿が、体外に排出されるときに通る管を何というか。}{尿道}

\Instruction{腎臓と汗腺のはたらきに関する次の問いに答えなさい。}
\QQ{\RefBlank{kidney}が余分な水分や塩分を取り除くことで、血液の何を一定に保つことができるか。}{成分や濃度}
\QQ{尿素などの不要な物質や、水分、塩分の一部を汗として排出する器官を何というか。}{汗腺}
\QQ{汗の成分は、何の成分とほぼ同じか。}{尿}
\QQ{汗は尿と比べて、水分が多いか少ないか。}{多い}
\QQ{汗は尿と比べて、濃度が濃いかうすいか。}{うすい}

\Instruction{肝臓のつくりとはたらきに関する次の問いに答えなさい。}
\QQ{体の中で最大の臓器を\NamedBlank{liver}という。}{肝臓}
\QQ{\RefBlank{liver}は、横隔膜の下にある腹部のどのあたりにあるか。}{右上部}
\QQ{\RefBlank{liver}でつくられ、胆のうにたくわえられた後、小腸に送り出される液体を何というか。}{胆汁}
\QQ{小腸で吸収された養分の一部をつくり変えたり、たくわえたりする器官は何か。}{肝臓}
\QQ{食物中の有害な物質を、無害な物質に変える器官は何か。}{肝臓}
